\chapter{Computation as Physics}
\label{sec:computation-as-physics}

Everyone says "information is physical." Cute slogan. Wrong level.

The claim in CΩMPUTER is stronger and less polite:

\begin{quote}
\textbf{Physics is what certain computations \emph{look like} from the inside.}\\
Conserved quantities, fields, particles, and thermodynamic laws are emergent
properties of a particular region of MRMW under a particular choice of coarse-graining.
\end{quote}

This chapter is where we stop treating that as a motivational poster and actually
spell out the mapping.

We are not "simulating physics." We are identifying physics with
\emph{statistics and invariants of rewrite} on a particular class of WARP graphs.
The universe isn’t running on a computer; the universe is what it feels like
to be a very specific CΩMPUTER instance.

\subsubsection*{25.1  Substrate First: WARP + DPO as "World"}

We start from the ontology built in Parts I–III:

\begin{itemize}
  \item A \textbf{WARP Graph} (WARP) is the underlying state object.
        Nodes and edges can themselves point to graphs, rules, or higher-order
        structures. "State" is not a flat configuration; it's a tower of graphs
        describing both objects and the laws that act on them.

  \item \textbf{Double-pushout rewrite} (DPO) is the evolution law:
        each step replaces a local pattern $L$ with a pattern $R$, glued into
        the ambient graph via an interface $K$:
        \[
          L \xleftarrow{i_L} K \xrightarrow{i_R} R
        \]
        plus a match $m : L \to G$ into the current world-graph $G$.

  \item \textbf{MRMW} is the phase space of all such evolutions:
        all $(G, \mathcal{R})$ pairs (graphs plus rule sets) and all their
        possible rewrite histories. Worldlines are particular rewrite histories,
        i.e., paths through MRMW.
\end{itemize}

This is not “an encoding.” This \emph{is} the substrate.

The first move in this chapter is: given this substrate, what has to be true
for an embedded observer to recognize something we would call
"physics" rather than arbitrary combinatorics?

\subsubsection*{25.2  What Counts as "Physics"?}

From within a universe, nobody sees DPO diagrams. They see:

\begin{itemize}
  \item regularities (the same experiment gives the same result),
  \item constraints (some configurations never occur),
  \item and scalable summaries (you don’t need to track every bit of sand to
        predict a sandpile).
\end{itemize}

Formally, that means:

\begin{enumerate}
  \item There exists a \textbf{coarse-graining} $C$ from the raw WARP state $G$
        to macrostates $X = C(G)$ such that the induced dynamics on $X$ is:
        \begin{itemize}
          \item approximately Markovian on relevant time scales,
          \item approximately local (future $X'$ mostly depends on a neighborhood of $X$),
          \item and stable under small perturbations of microstate.
        \end{itemize}

  \item There exists a family of \textbf{invariants} under the DPO dynamics,
        i.e., quantities $I(G)$ such that $G \to G'$ implies $I(G) = I(G')$,
        which appear to embedded agents as "conserved’’ (mass, charge, etc.).

  \item There exists a notion of \textbf{causal structure} induced by rewrite:
        a partial order on events where some pairs are spacelike (independent),
        some timelike (influence possible), and maximal influence speed is bounded.
\end{enumerate}

Whenever these conditions hold, embedded agents will end up inventing something
indistinguishable from "physics," whether or not they know WARP graphs or DPO exist.

The rest of the chapter is just spelling out those bullets in WARP language.

\subsubsection*{25.3  Conservation Laws as Rewrite Invariants}

Take a DPO rule:
\[
  L \xleftarrow{i_L} K \xrightarrow{i_R} R
\]

We define a \textbf{quantity} $Q$ as a function from graphs to some value space
(e.g., $\mathbb{Z}$, $\mathbb{R}^n$, an abelian group, etc.).
We say $Q$ is \textbf{locally conserved} by a rule if:
\[
  Q(L) = Q(R)
\]
under an appropriate notion of sum over components, and \textbf{globally conserved}
if this local equality holds for all matches of all rules in $\mathcal{R}$.

Physically:

\begin{itemize}
  \item "Charge" is any additive $Q$ which counts certain labeled substructures.
  \item "Particle number" is $Q$ counting connected components in some layer.
  \item "Topological charge" is $Q$ depending only on global features
        (e.g., winding number-like quantities over the graph embedding).
\end{itemize}

Important: none of these are declared by fiat. They are discovered as invariants
of the rule system, just like conservation laws in Noether-style physics are
discovered from symmetry.

The move here is: we treat \emph{every} invariant $Q$ induced by a rule set
as a candidate "physical quantity" for the embedded universe. Most will be
uninteresting; the physically relevant ones are those that survive coarse-graining
and show up in macroscopic behavior.

\subsubsection*{25.4  Symmetry as Rule Equivariance}

In usual physics, symmetry is group action on a state space that commutes with dynamics.

In CΩMPUTER:

\begin{itemize}
  \item A \textbf{symmetry} is a group action $S: G \mapsto S(G)$ on graphs
        (and optionally on rules) such that applying $S$ then evolving by DPO
        is equivalent to evolving then applying $S$.

  \item Formally, for each rewrite $G \to_{\rho} G'$, there is a corresponding
        rewrite $S(G) \to_{S(\rho)} S(G')$, where $S(\rho)$ is the pushed-forward rule.
\end{itemize}

If $S$ is a spatial translation, you get homogeneity.  
If $S$ is a rotation, you get isotropy.  
If $S$ is time-shift, you get energy conservation via the $Q$ induced by that symmetry.

This lets us lift the usual physics story:

\begin{center}
\renewcommand{\arraystretch}{1.2}
\begin{tabularx}{\linewidth}{>{\bfseries}l X X}
\toprule
Physics & CΩMPUTER & Emergent meaning\\
\midrule
Spatial translation invariance & Graph automorphisms preserving pattern frequencies & No preferred location \\
Rotational invariance & Automorphisms of embedding / adjacency patterns & No preferred direction \\
Time invariance & Rule set constant along worldline & Energy-like invariant exists \\
Gauge symmetry & Local relabeling of subgraphs & Redundant description of “fields” \\
\bottomrule
\end{tabularx}
\end{center}

We’re not copying physics; physics drops out as the special case where the
symmetry group is large and the coarse-graining is nice.

\subsubsection*{25.5  Locality, Causality, and Light Cones}

DPO rewrite is \emph{by construction} local: you only touch the match of $L$
and its glued neighborhood. But that doesn’t automatically give you nice spacetime.

We define:

\begin{itemize}
  \item A \textbf{causal edge} from event $e$ to event $e'$ if the application
        of some rule at $e$ changes the part of the graph that is matched at $e'$.

  \item A \textbf{causal graph} whose nodes are rewrite events and edges are
        these dependencies.

  \item A \textbf{light cone} of an event $e$ as the subgraph induced by
        all events reachable from $e$ in the causal graph (future cone) and
        all events that can reach $e$ (past cone).
\end{itemize}

We say a rule set is \textbf{bounded influence} if there exists a finite $v$
such that the radius of the affected region after $t$ steps grows at most like $vt$.

For an embedded observer, this bounded influence appears as:

\begin{itemize}
  \item a maximum signal speed (a "speed of light"),
  \item a separation between timelike and spacelike separation,
  \item and causality constraints (no signaling outside light cones).
\end{itemize}

If you want relativity, you need more: you want the causal structure to be
robust under coarse-graining and compatible with a family of emergent frames.
But the skeleton is already here: causal graph first, metric structure second.

\subsubsection*{25.6  Thermodynamics as Combinatorics on MRMW}

Now the spicy part: thermodynamics without waving your hands about "ensembles"
as something outside the universe.

We define:

\begin{itemize}
  \item A \textbf{microstate} as a full WARP graph $G$.
  \item A \textbf{macrostate} as an equivalence class $[G]$ under some coarse-graining $C$.
  \item The \textbf{entropy} of a macrostate $X$ as:
    \[
      S(X) = \log \left|\{ G \mid C(G) = X\} \right|
    \]
    where the log base sets units (pick your poison).
\end{itemize}

Given a rule set and a coarse-graining, typical behavior is:

\begin{enumerate}
  \item Under evolution, the system explores microstates consistent with
        its invariants (conserved $Q$’s).
  \item The overwhelming majority of microstates correspond to high-entropy
        macrostates.
  \item Embedded agents, having access only to macroscopic observables, describe
        this as "systems evolve toward equilibrium."
\end{enumerate}

Temperature, free energy, etc. emerge as parameters describing how rewrite
trajectories distribute over macrostates:

\begin{itemize}
  \item \textbf{Temperature} is effectively a Lagrange multiplier controlling the
        trade-off between energy-like invariants and entropy in the reachable region.
  \item \textbf{Free energy} is the "computational slack":
    \[
      F = \langle E \rangle - TS
    \]
    where $E$ is some energy-like $Q$ and $S$ is macrostate entropy.
\end{itemize}

This is not metaphor. It’s literally counting the ways the graph can be while
keeping some invariants fixed.

\subsubsection*{25.7  Landauer in WARP Clothing}

Landauer’s principle says: erasing one bit of information has a minimum
thermodynamic cost.

In CΩMPUTER, this becomes:

\begin{itemize}
  \item A \textbf{logically irreversible} rewrite is one that maps many distinct
        microstates to the same microstate class, i.e., the preimage of the
        rule application is larger than the image.

  \item Implementing such a rewrite inside a larger DPO system necessarily
        pushes entropy somewhere else in the graph: you can’t compress the
        preimage set without expanding the state space of some environment.

  \item The energy cost is just what an embedded observer calls the work needed
        to perform those compensating rewrites against their preferred gradients.
\end{itemize}

We recover the same intuition: information destruction is not free; it’s just
graph structure being scrambled elsewhere.

This perspective is particularly useful for Part V, where we’ll design runtimes
that track logical and physical irreversibility separately.

\subsubsection*{25.8  Time’s Arrow as Coarse-Graining Bias}

Micro DPO rules can be perfectly reversible: every step has an unambiguous inverse.
Yet embedded observers still see an arrow of time.

In WARP language:

\begin{itemize}
  \item The \textbf{microscopic dynamics} can be symmetric: for every allowed
        rewrite $G \to G'$ there is an inverse rewrite $G' \to G$.

  \item The \textbf{coarse-graining} $C$ is almost never symmetric: many microstates
        map to the same macrostate, and the typical flows in macrospace go from
        low-entropy to high-entropy regions.

  \item "Irreversibility" is a property of $(\mathcal{R}, C)$, not of $\mathcal{R}$ alone.
\end{itemize}

In other words, the arrow of time is a perspectival artifact of the particular
macro-variables the embedded agents choose to track. Once they compress away
microstate identity, forward and backward trajectories in macrospace look very
different.

This lines up with standard statistical mechanics, but we now have a concrete
computational substrate beneath it: no appeal to vague ensembles, just explicit
counts of WARP states.

\subsubsection*{25.9  Mapping Table: Physics $\leftrightarrow$ CΩMPUTER}

We can now summarize the core identifications:

\begin{center}
\begin{tabular}{ll}
\toprule
\textbf{Physics Concept} & \textbf{CΩMPUTER / WARP Concept} \\
\midrule
State of the world & WARP $G$ (possibly with layered structure) \\
Law of motion & DPO rule set $\mathcal{R}$ \\
Event & Single DPO rewrite application \\
Spacetime point & Node in causal graph of rewrite events \\
Worldline & Path in causal graph / MRMW trajectory \\
Conserved quantity & Invariant $Q(G)$ under all rules \\
Symmetry & Group action commuting with DPO dynamics \\
Locality & Bounded rewrite match radius; sparse causal graph \\
Light cone & Reachable set in causal graph from event \\
Entropy & $\log$ of microstate count per macrostate \\
Temperature & Lagrange multiplier on invariants in MRMW measures \\
Free energy & Trade-off between invariant $Q$ and entropy \\
Arrow of time & Asymmetry induced by coarse-graining $C$ \\
\bottomrule
\end{tabular}
\end{center}

This table is not decorative. It’s the checklist for deciding whether a given
WARP + rule set deserves to be taken seriously as a "physical" universe.

\subsubsection*{25.10  Why This Matters (and What Comes Next)}

So what did we actually do in this chapter?

\begin{itemize}
  \item We \emph{reversed the metaphor}: physics is not something we simulate;
        it is a structured pattern inside MRMW.

  \item We grounded standard physical notions (conservation, symmetry, light cones,
        entropy, time’s arrow) directly in WARP + DPO.

  \item We set up a concrete program: given a candidate rule system,
        you can now ask, in a disciplined way, "Does this generate physics-like
        behavior under some sane coarse-graining?"
\end{itemize}

The next chapter, \emph{Physics as Computation}, will push in the opposite direction:

\begin{quote}
Given a traditional physical theory (say, a Lagrangian or a differential equation),
how do we \emph{compile} it into a WARP + DPO rule set such that the mapping
outlined here becomes exact rather than hand-wavy?
\end{quote}

That’s where we start taking specific chunks of real-world physics and
turning them into concrete CΩMPUTER universes, one rule system at a time.
